\chapter{细胞合作形成身体}

\begin{kaichang}
回想一下你上一次受的小伤：也许是削铅笔划破了手指，也许是膝盖磕掉了一小块皮。当时流血、刺痛，你贴了创可贴，几天后边缘长出嫩粉色的新皮，再过些日子，结痂翘起、脱落，皮肤几乎看不出痕迹。

现在请你认真回答一个问题：\textbf{那道伤口是怎么“知道”自己缺了一块，又是怎么知道该补多少、什么时候该停工的？}没有工程师在场指挥，没有图纸贴在伤口边上。补少了是个洞，补多了是个鼓包——可大多数时候，修补恰恰停在刚刚好的位置。

先写下你的猜测。这一章结束时我们回到这道伤口。答案会牵扯出一个比你身体大得多的问题：几万亿个细胞，凭什么愿意——不对，凭什么\emph{能够}——合作成一个人？
\end{kaichang}

\section{现象层：一道伤口的慢动作直播}

先把分子和细胞都放下，只看你能观察到的。

受伤后几分钟内，血会变稠、凝住——这就是止血。接下来一两天，伤口周围发红、发热、微微肿痛，摸上去比旁边烫。再往后，伤口底下长出颗粒状的嫩红色组织，医生叫它“肉芽”，它像填充材料一样把坑垫平。与此同时，伤口边缘的皮肤开始向中间“爬”，一点点把表面盖上。最后结痂脱落，新皮肤颜色偏浅，如果伤得深，会留下一道疤，而这道疤在此后几个月里还会慢慢变软、变淡。

还有一些你平时不会在意的现象，其实是同一出戏的日常版：洗澡搓下来的“泥”里有大量脱落的表皮细胞；剪掉的指甲过几周又长出来；晒伤后皮肤成片脱皮；口腔溃疡疼几天就自己好了。这些现象共同说明一件事：\textbf{你的身体不是建好之后一成不变的雕像，而是一个不停拆旧补新的工地}——受伤只是把工地的活儿加快了、集中了，让我们有机会看见它。

\begin{shiyan}
\textbf{给你的指甲装一把“游标卡尺”}

材料：一支可水洗的细头彩笔、一把毫米刻度尺、一个日历（或手机备忘）。

步骤：选一只手的拇指指甲，在靠近甲根（月牙旁边）的皮肤与指甲交界处，用彩笔在指甲上画一条细横线作为基准线，记下日期。之后每隔七天，用尺子量一次基准线到甲根皮肤的距离，记录下来，连续观察三到四周。

观察点：每周长出的距离是不是差不多？算出平均每天长多少毫米。再想想：基准线本身为什么不会被“拉宽”？（提示：已经长出来的指甲还是活细胞吗？）

安全提示：只画在完整的指甲上，皮肤破损处不要画；不要用小刀在指甲上刻记号。指甲平均每月长约 \ud{3}{mm}，不同人、不同手指差异很大，你的数据偏快偏慢都正常。
\end{shiyan}

\section{粒子层：谁赶来修补这座城市}

把镜头推进到伤口边缘，那里热闹得像事故现场。

最先到场的是\textbf{血小板}——血液中的一种细胞碎片，它们碰到破损的血管壁就被激活，变形、聚集，像一袋袋速干水泥一样堵住缺口；同时血液里的凝血因子启动连锁反应，把可溶的纤维蛋白原切成纤维蛋白，织成一张网，把血细胞兜住——这就是血痂的主要成分（凝血这条酶促级联的原理，正是第 49 章讲的酶促反应的活用）。

紧接着，几类“常驻维修队”开始干活。免疫细胞清理细菌和坏死组织——这一支队伍的识别、反应和记忆极其精彩，我们留到第 61 章专门讲。\textbf{成纤维细胞}则像装修工人，一边分泌胶原蛋白和其他基质材料把坑填满，一边拉着伤口边缘往里收。伤口边缘的\textbf{上皮细胞}改变形状，松开彼此之间的连接，像摊开的手掌一样向中间滑行、铺开，边滑行边分裂，直到覆盖整个创面，然后重新叠成多层、角化，恢复皮肤的屏障功能。与此同时，新的毛细血管像水管一样长进肉芽组织，给工地送料送氧。

请你注意一个关键事实：血小板、上皮细胞、成纤维细胞、血管内皮细胞，\textbf{长相完全不同，干的活完全不同，可它们全都是你身上的细胞}。第 57 章讲过，每个细胞内部都有细胞核、线粒体、细胞骨架这一整套设备；第 58 章讲过，细胞靠信号分子和受体“听见”外界的命令。但这两章留下了一个悬而未决的问题：既然大家装着同一套 DNA，凭什么你长成上皮细胞，我长成肌肉细胞？

\section{规则层：同一本说明书，读不同的页}

你身上几乎每个细胞核里的 DNA 都一模一样——都来自同一个受精卵的一分为二、二分为四。差别不在“说明书”的内容，而在\textbf{每一类细胞只打开其中一部分页面}。这个过程叫\textbf{细胞分化}：一个细胞的后代在基因表达的选择上逐渐分道扬镳，走上不同的“职业道路”。

用人话说：DNA 是一座有几万个菜谱的图书馆。上皮细胞只借“生产角蛋白、紧密排列”那一批书，肌肉细胞只借“生产肌动蛋白、肌球蛋白（第 48 章讲过的收缩机器）”那一批，神经细胞借的是“离子通道、神经递质”那一批。哪些书能借出来，由一类叫转录因子的蛋白质控制——它们识别 DNA 上特定的开关序列，决定附近的基因是开工还是沉默。而转录因子本身又听命于细胞收到的信号：邻居是谁、基质软硬、激素浓度（第 58 章的信号通路在这里接上头）。

分化通常是一步一步定下来的，像不断收窄的路口：

\begin{center}
\begin{tikzpicture}[scale=0.9, every node/.style={font=\small}]
\node[draw, rounded corners] (z) at (0,0) {受精卵};
\node[draw, rounded corners] (s) at (0,-1.3) {各类祖细胞};
\node[draw, rounded corners] (a) at (-4,-2.6) {上皮细胞};
\node[draw, rounded corners] (b) at (-1.3,-2.6) {肌肉细胞};
\node[draw, rounded corners] (c) at (1.3,-2.6) {神经细胞};
\node[draw, rounded corners] (d) at (4,-2.6) {血细胞等};
\draw[->] (z) -- (s);
\draw[->] (s) -- (a);
\draw[->] (s) -- (b);
\draw[->] (s) -- (c);
\draw[->] (s) -- (d);
\end{tikzpicture}
\end{center}

这是帮助思考的表示法：真实的分化不是一次分叉到位，而是经过许多中间状态逐步收窄的。

\begin{zhengju}
“分化只是把用不到的页面合上了，还是把整页撕掉了？”——这是二十世纪中期发育生物学的头号争论。如果分化伴随基因丢失，那么一个彻底分化的细胞就永远回不了头；如果只是表达被关闭，理论上细胞核里还保存着完整的说明书。

1962 年，英国科学家格登做了一个干脆的判决实验：他把非洲爪蟾小肠上皮细胞（已经彻底分化、专门负责消化吸收的细胞）的细胞核取出来，移植到一个被去掉细胞核的卵细胞里。结果，其中一些“借核重生”的卵竟然发育成了正常的蝌蚪，甚至长成可育的成蛙。要知道，在那之前布里格斯和金用早期胚胎细胞核做类似移植，结果时好时坏，很多人据此认为分化一旦走远就不可逆。格登的结果说明：\textbf{肠细胞核里一个字都没少}，整套建造青蛙的指令都还在，只是大部分被关掉了；卵细胞的细胞质里有办法把它们“重置”。

这个故事还没完。2006 年，山中伸弥更进了一步：不用卵细胞，只用四种转录因子，就能把成年小鼠的皮肤细胞直接“倒带”回类似胚胎干细胞的状态（诱导多能干细胞，iPS 细胞）。格登与山中伸弥因此分享了 2012 年诺贝尔生理学或医学奖。请注意这个证据链的形状：先有一个可被推翻的问题，再有判决性实验，中间经历怀疑和修正，最后连“怎么做到的”都被一步步拆解成具体分子。这才是“科学家怎样知道”的真实样子。
\end{zhengju}

\section{四种“建筑材料”：组织}

分化的结果，是身体出现了几百种细胞。把结构相似、功能相关的细胞连同它们分泌的材料放在一起，称为\textbf{组织}。传统上把动物组织分成四大类——记住，这是人为的归纳，真实的细胞类型远比四类丰富。

\begin{itemize}[itemsep=2pt]
\item \textbf{上皮组织}：身体的“外墙和内衬”。皮肤表皮、消化道内壁、肺泡壁都是它。上皮细胞紧密排列、彼此锁死，一面朝外、一面朝内，负责屏障、吸收和分泌。它处在磨损第一线，所以更新也最快。
\item \textbf{结缔组织}：填充、连接和支撑。血液、脂肪、软骨、骨、肌腱都属于这一类。它的特点是细胞少、“细胞间的东西”多——骨头之所以硬，主要不是骨细胞硬，而是细胞分泌的材料硬。
\item \textbf{肌肉组织}：专职收缩。骨骼肌听你指挥，心肌终生不停，平滑肌默默推着食物前进、调节血管粗细。收缩的分子原理（肌动蛋白与肌球蛋白的相对滑动）在第 48 章已经见过。
\item \textbf{神经组织}：信息的高速公路。神经元用长长的突起和电化学信号把全身连成一张网；第 58 章讲过的神经递质，就是这张网传递消息的“包裹”。
\end{itemize}

这里要特别为一个常被忽视的角色正名：\textbf{细胞外基质}。教科书的细胞图常把细胞画成挤在一起的球，球与球之间一片空白。真实的组织里，细胞浸泡在自己分泌的“凝胶和钢索”里：胶原纤维提供抗拉强度（第 48 章讲过这种蛋白质的三股螺旋），弹性蛋白提供回弹，蛋白聚糖吸住水分提供抗压的“水垫”。皮肤有韧性、软骨能缓冲、肌腱拉不断，靠的主要是这些细胞外材料。

更重要的是，基质\emph{不是}无生命的填充物。细胞膜上有一类叫整合素的蛋白质，一头抓住细胞外的基质，一头连着细胞内的细胞骨架（第 57 章）——细胞借此“摸”出环境的软硬和拉力，并据此调整自己的分裂、移动和分化。干细胞放在软基质上容易走向神经方向，放在硬基质上容易走向骨的方向。伤口处的成纤维细胞正是抓着胶原纤维收缩，才把伤口拉拢的。基质的成分和力学性质还会随年龄、疾病改变：皮肤老化松弛、动脉粥样硬化斑块变硬，都有基质变化的份。

\section{合作的条约：更新、储备与退出}

一个由几万亿成员组成的社会，最怕两件事：没人接班，和赖着不走。多细胞身体用两条对应的规则解决。

第一条是\textbf{持续更新}。你每天脱落的表皮细胞数以十亿计，肠道内壁每三到五天就整体换一遍，红细胞的寿命约 120 天，到期就被脾脏回收。补员从哪来？来自各组织里的\textbf{干细胞}：它们分裂时，一个子细胞留下来继续当干细胞，另一个走向分化，成为接班的新鲜劳动力。骨髓里的造血干细胞源源不断地制造红细胞、白细胞和血小板；小肠内壁凹陷处的干细胞几天就能把肠壁表面全部刷新；皮肤基底层也驻守着干细胞的储备。

第二条是\textbf{程序性细胞死亡}，学名叫凋亡。听起来可怕，可它恰恰是合作得以维持的关键。细胞内部的损伤探测器（比如第 48 章提到过的蛋白质质量监控的升级版）发现 DNA 损伤修不好、或者收到“你多余了”的信号时，细胞会启动一套自毁程序：DNA 被切成小段，细胞皱缩成小块，被邻居或免疫细胞安静地吃掉，不泄漏、不发炎。这和意外死亡（坏死）完全不同——坏死是细胞被撑破，内容物泼洒一地，引发炎症。

凋亡少一次，后果可能是癌症；凋亡多过头，后果可能是组织萎缩。它甚至参与建造：胚胎发育时，你的手指之间最初是有“蹼”连着的，指间细胞按计划凋亡，五指才分开；蝌蚪变态成蛙时尾巴消失，靠的也是凋亡，而不是尾巴“掉了”。

\begin{qiaoliang}
你的身体里大约有多少细胞在“交接班”？做一个数量级估算。科学家估计成年人体细胞总数约为 $3.7\times 10^{13}$ 个（三十七万亿，注意这本身是粗略估计，不同算法能差出好几倍）。其中更新最猛的是血液和肠道：仅红细胞每天就要新生约 $2\times 10^{11}$ 个。换算一下：
\[
\frac{2\times 10^{11}}{24\times 3600\ \text{秒}} \approx 2\times 10^{6}\ \text{个/秒}
\]
也就是说，此刻你读这句话的一秒钟里，骨髓里大约有\textbf{两百万个红细胞}诞生，同时脾脏里大约有两百万个到期的红细胞被拆解回收。把血液、肠道、皮肤加在一起，你每天更替的细胞总数约 $3\times 10^{11}$ 个，占细胞总量的比例不到百分之一——听上去不多，一年累计下来，相当于把全身细胞按重量换掉了相当一部分。衰老，在某种程度上就是这套“交接班系统”的精度在慢慢下降。
\end{qiaoliang}

程序性死亡的精确程度，被一条小虫暴露得一清二楚。

\begin{zhengju}
秀丽隐杆线虫是一种身长约 \ud{1}{mm} 的透明小虫，它的成虫全身只有 959 个体细胞——少到可以挨个编号、逐个追踪。二十世纪七八十年代，萨尔斯顿等人借着它的透明身体，在显微镜下把一个受精卵的每一次分裂、每一个细胞的去向全部记录下来，画出了完整的“细胞家谱”。结果发现一个惊人的规律：发育中总共产生 1090 个细胞，其中不多不少恰好 131 个，在固定的时刻、固定的位置死掉，而且每一只线虫都一样。

死得这么整齐划一，不可能是事故。霍维茨接着找到几个关键基因（ced-3、ced-4、ced-9 等）：ced-3 失活的线虫，那 131 个细胞一个都不死，全部活下来——说明死亡是细胞主动执行的程序，由基因编码。后来人们发现，这些基因在包括人在内的动物里都有亲戚，它们出故障与人类癌症、神经退行性疾病密切相关。2002 年，布伦纳、霍维茨、萨尔斯顿因这项工作获得诺贝尔奖。一条透明小虫，把“细胞自杀也是发育的一部分”这个反直觉的结论，钉成了铁的事实。
\end{zhengju}

\section{当合作破裂：癌症是一种演化}

现在可以正面回答“凭什么合作”了。答案是：\textbf{合作不是细胞的品德，而是演化的结果}。

多细胞生命在地球上独立起源了许多次。今天与动物亲缘最近的单细胞生物——领鞭毛虫——已经拥有不少细胞间粘连和通讯的蛋白质零件，说明“合作的零件箱”在多细胞出现之前就在组装。一旦一群细胞共享同一套基因组、并且只有当整个群体活下去时自己的基因才能传下去，自然选择就会偏爱那些限制分裂、服从信号、必要时自我了断的细胞。你的身体，本质上是几十亿年演化筛选出的合作制度。

但这个制度时刻面临“作弊者”：某个细胞碰巧积累了让分裂失控、让凋亡失灵的突变组合，它就不再遵守条约，只顾自己增殖——这就是肿瘤。第 58 章讲过癌细胞怎样篡改生长信号，这里换一个视角看：\textbf{肿瘤内部就是一个小小的演化实验场}。细胞分裂时不断产生新突变，身体的环境（免疫攻击、缺氧、药物）充当选择压力，最能逃避免疫、最耐药的细胞后代存活扩张。化疗有时先有效后失效，原因和细菌产生耐药性（第 62 章会讲）是同一套演化逻辑：治疗杀死敏感细胞，等于替耐药突变清了场。

演化视角还带来一个著名的谜题——佩托悖论：大象的细胞数量约是人类的一百倍，寿命也不短，理论上每次分裂都是一次出错机会，大象的癌症风险应该高得多，可实际上并不如此。研究提示大象的基因组里有几十份抑癌基因 p53 的拷贝（人类只有两份），DNA 损伤一出问题就坚决凋亡。自然选择在大体型动物身上加码了“警察系统”。细节仍在研究中，但方向很清楚：抗癌机制本身，也是演化的产物。

\section{边界层：这幅图景在哪里会失真}

本章的叙述有几个简化，你应当知道它们什么时候会失灵。

第一，“上皮、结缔、肌肉、神经”四大组织是两百年前显微镜时代的分类。现代单细胞测序发现，真实组织里的细胞状态更像连续的光谱而不是离散的抽屉：同一种“巨噬细胞”在不同器官里状态差异很大，还不断随环境变化。四大类仍然是好用的地图，但别把它当成格子铺。

第二，“分化是单行道”也在被改写。格登的实验和 iPS 技术说明细胞命运可以被重置；在肝脏等器官里，已分化的细胞在受损时能“转回头”补充其他类型。不过，不要把实验室里的可塑性当成日常生活的常态——人体内绝大多数分化是稳定的，疤痕之所以是疤，恰恰因为修复并不总能完美重来。

第三，本章讲“伤口愈合”说的是皮肤小伤口的标准剧本。严重烧伤、慢性伤口（如糖尿病足）会卡在某一步走不下去；脊髓和大脑的神经元几乎不能再生。修复能力在不同组织、不同物种间差异巨大——蝾螈能长出整条腿，我们只能长出一道疤，原因至今没有被完全解开。

\begin{wuqu}
“细胞分化就是把用不到的基因扔掉了，所以每个细胞只装自己需要的那部分 DNA。”——这是最流行的错误图景之一。反例一：把胡萝卜根的一块韧皮部组织打散，单个分化的细胞在培养基里能重新长成一株完整的、会开花的胡萝卜（植物组织培养，如今已是农业育苗的常规技术）。反例二：格登的核移植实验（本章证据故事）和克隆羊多莉——一个乳腺细胞的核，就足以指导发育出一只完整的羊。如果基因真的被扔掉了，这些都绝无可能。真相是：\textbf{分化的差别在基因表达的开关状态，而不在基因的有无}。这个区别不是咬文嚼字——克隆技术、再生医学的全部希望，都建立在“说明书还在”这一点上。
\end{wuqu}

\section{回到那道伤口}

现在，把开场的问题拆开重答。

伤口怎么“知道”缺了一块？——完整的上皮层里，细胞彼此紧密接触，接触本身就抑制分裂；缺口一出现，边缘细胞失去了侧面的接触和拉扯，加上血小板释放的生长因子（第 58 章的信号分子），等于同时收到“解除刹车”和“踩油门”两条命令，于是开始分裂、向中间迁移。

它怎么知道该补多少？——当两侧的上皮细胞在伤口中央相遇，细胞间的接触重新建立，刹车信号恢复，分裂和迁移随之停止。坑里的肉芽组织随后被成纤维细胞逐步改建成更结实的胶原，多余的新生血管凋亡退去。

它怎么知道什么时候停工？——没有“知道”，只有反馈：命令在完成任务的那一刻自动消失。这套逻辑和家里的空调一模一样：不是空调“懂得”室温，而是传感器、控制器和执行器连成了一个闭环。下一章你会发现，你的体温、血糖、血压，全都靠这样的闭环维持。

至于那道疤：它是“快速抢修”而非“原样复刻”的遗迹——胶原纤维排得又密又乱，没有毛囊和汗腺。演化在这笔账里选择了速度：对野外的祖先来说，尽快封住伤口防感染，比长得好看重要得多。

\section{继续深挖：从修理身体到制造组织}

理解了细胞合作的规则，人类就开始学着插手。造血干细胞移植（骨髓移植）用健康供者的干细胞重建病人的整个血液系统，是治疗白血病的成熟手段——它之所以可行，正因为干细胞能长期自我更新并分化出所有血细胞类型。大面积烧伤的病人可以接受自体表皮培养移植：取一小块皮肤，在实验室把表皮细胞扩增成大片，再植回创面。科学家还在尝试用 3D 打印和支架材料（模拟细胞外基质）制造更复杂的组织。

但也要给你泼一盆冷水：市面上大量打着“干细胞”旗号的美容、抗衰项目，证据等级与上述正规疗法天差地别，多数缺乏严格的对照试验支持，有些甚至有诱发肿瘤的风险——毕竟，促进细胞分裂正是干细胞的本职，分裂失控的边界并不遥远。判断一个疗法，永远先问：证据来自哪里？剂量和风险说清楚了吗？

另一个前沿是再生：为什么蝾螈能再生心脏和脊髓而人不能？答案可能藏在细胞外基质、炎症反应和分化状态的差异里。如果哪一天科学家找到了让人体组织“安全地倒带”的办法，那将是本章所有规则的集大成应用。

而这一切——伤口的修补、血液的更新、温度的维持——都要求全身几万亿细胞在同一个化学环境里同步运行。血糖高了会怎样？体温低了会怎样？pH 偏一点会怎样？身体凭什么能把这些变量死死摁在狭窄的范围里？那是下一章的主题：稳态。

\begin{tiaozhan}
\textbf{从一个细胞到一个人，需要分裂多少次？}（跳过不影响主线）

假设身体约有 $3.7\times 10^{13}$ 个细胞，全部来自受精卵的一路分裂。如果每次分裂细胞数翻倍，需要分裂 $n$ 代使得 $2^n \approx 3.7\times 10^{13}$。两边取以 2 为底的对数：
\[
n = \log_2(3.7\times 10^{13}) \approx 45
\]
只要 45 代左右！指数增长的力量就在这里。但真实情况复杂得多：发育中大量细胞会按计划凋亡（比如那 131 个线虫细胞），不同组织的分裂次数差别巨大，而且像神经元这样的细胞在出生后基本停止分裂，脂肪细胞数量却会随体重变化。所以 45 只是一个粗糙的下限估算——它真正的价值是让你对“$10^{13}$ 个细胞从何而来”有一个数量级直觉，而不是一个精确答案。
\end{tiaozhan}

\section*{练一练}

\begin{sikaoti}
\item 画一幅“伤口物质流图”：方框代表血痂、肉芽组织、新上皮、疤痕，箭头标注每个阶段是谁（血小板、成纤维细胞、上皮细胞、血管）把什么材料送到了哪里。不必画细胞的样子，重在把流程连起来。
\item 上皮细胞和神经细胞的 DNA 几乎完全相同，功能却天差地别。用“图书馆只借出不同书架的书”这个比喻，解释分化；再用格登的核移植实验说明：为什么“书被合上”和“书被撕掉”这两种假说，实验结果只支持前者？
\item 为什么伤口边缘的上皮细胞迁移到中央相遇后就会停下来？把“接触抑制”改写成一个负反馈回路（刺激→响应→刺激消失），并再举一个生活中同类的反馈例子。
\item 凋亡既参与雕刻手指，也参与清除癌前细胞。有人说“细胞死亡总是坏事，应该用药阻止”，请用本章的论据反驳他；再说明为什么“凋亡越多越好”同样站不住脚。
\item 用正文的数据估算：一个红细胞寿命约 120 天，每天新生约 $2\times 10^{11}$ 个，那么你体内红细胞总数大约是多少？（用乘法即可，结果用科学记数法表示。）
\item 佩托悖论说：细胞越多、寿命越长的动物，按直觉癌症应该越多，但大象并非如此。从“演化会塑造合作制度”的角度，提出一种解释，并说说可以收集什么证据来检验你的解释。
\end{sikaoti}
